% ============================================================
% 经典洞见（杂志页码 54–59）
% 「村BA」专题 + 技法指南 + 考场范本
% ============================================================
\colchapter{经典洞见}

\section*{「村BA」火爆：有温度的乡土文化，接地气的乡村振兴}
\auth{编写者：广东省佛山市禅城实验高中 周星明}

2022年夏天，贵州省台江县台盘村的一场乡村篮球赛火了。因为场地设在村里、参赛的是村民、观众席里三层外三层，火热的氛围堪比职业联赛，网友便戏称它为“村BA”。

“村BA”为什么这么火？归根结底，在于一个“土”字。比赛没有专业的场馆，就设在村口的篮球场；球员没有职业的装备，大多是放下锄头的村民；观众没有整齐的座位，就围在球场边或站或蹲。但正是这种“土”，让比赛充满了生活的烟火气：中场休息时，村民端上苗家酸汤鱼；啦啦队上场，跳的是苗族的芦笙舞；奖品不是奖杯，而是台江的农产品。

“村BA”的背后，是深厚的群众基础。篮球在台盘村有着几十年的传统，每逢节日，村里都要举办篮球赛，这已经成为村民生活方式的一部分。当体育运动真正融入乡村生活，它就不再只是“看个热闹”，而成为一种凝聚人心、传承文化的力量。

更重要的是，“村BA”让我们看到了乡村振兴的另一种可能：不只是一味地“输血”搞建设，而是激活乡村自身的文化活力，让村民成为乡村文化的主角。有温度的乡土文化，就是接地气的乡村振兴。

\begin{quoteBox}{文化知点}
台盘村位于贵州省黔东南苗族侗族自治州台江县，是“苗族刺绣之乡”。当地苗族群众有在“吃新节”等传统节日举办篮球赛的习俗，至今已延续数十年。“村BA”走红后，带动了当地农产品销售与乡村旅游发展，成为观察乡村振兴的一个窗口。
\end{quoteBox}

\clearpage
\section*{“村BA”：“从白天打到黑夜”，为什么这么燃？}
\auth{文/晶微}

最近，贵州台盘村的篮球赛火了。没有职业球员，没有商业赞助，一场乡村篮球赛，却让全网沸腾——“从白天打到黑夜”，围观的村民里三层外三层，直播间里数十万人同时在线。

“村BA”为什么这么燃？燃就燃在一个“真”字。球是真的，汗水是真的，欢呼是真的。参赛的球员，白天在地里干活，晚上在球场上奔跑；围观的观众，拖家带口搬着小板凳来看球，看到精彩处，嗓子都喊哑了。没有剧本，没有彩排，有的只是最纯粹的竞技热情。

燃就燃在一个“土”字。奖品是牛、羊、猪，啦啦队是村里的芦笙队，解说员讲着地道的贵州方言。这种带着泥土气息的“土味”，恰恰是“村BA”最大的魅力所在。它让篮球从职业赛场走回了田间地头，让每个人都能感受到体育最本真的快乐。

燃就燃在一个“根”字。台盘村打篮球的历史，可以追溯到几十年前。每逢“吃新节”，村里都要组织篮球赛，这是祖祖辈辈传下来的传统。当体育扎根于乡土，它就不再是悬浮的活动，而是融进了乡村的血脉里。可以说，“村BA”燃的，不只是篮球，更是乡土中国的生命力。

“村BA”的火爆，也为乡村文化振兴提供了启示：不必盲目追逐“高大上”，尊重当地传统、激活内生动力，让群众当主角，乡土文化就能焕发出最蓬勃的生机。

\clearpage
\section*{守住“村BA”的乡土味}
\auth{文/常理}

“村BA”火了之后，如何让这份热度延续下去，成为许多人思考的问题。有人建议引入商业赞助，有人提议扩建场馆，还有人想把“村BA”办成全国性的赛事。这些想法固然有可取之处，但在我看来，无论“村BA”走得多远，都不能丢了它的“乡土味”。

乡土味，是“村BA”的根。村民们喜欢的，不是职业联赛式的炫技，而是球场上那份亲切和热闹。如果比赛被商业过度包装，球员被职业选手取代，观众席变成整齐划一的看台，那么“村BA”也就失去了它最动人的东西。守住乡土味，就是要守住它的质朴与本真。

守住乡土味，也不等于故步自封。“村BA”可以借助网络传播扩大影响力，可以与乡村旅游、农产品销售相结合，让篮球赛成为展示乡村的窗口。但所有这些“加法”，都应该服务于乡村和村民本身，而不是本末倒置。

说到底，“村BA”属于乡村，属于村民。守住乡土味，就是尊重乡村自身的发展逻辑；让村民在文化活动中收获快乐与自信，才是“村BA”最大的价值。

\begin{quoteBox}{视野拓展}
台盘村“吃新节”是苗族同胞庆祝丰收的传统节日。节日期间，除了篮球赛，还有斗牛、芦笙舞、长桌宴等活动。“村BA”的走红，让更多人了解了这一传统节日，也带动了当地特色农产品——台江鲤鱼的销售。
\end{quoteBox}

\clearpage
\section*{“村BA”火热“出圈”，呼唤乡村文化自信}
\auth{文/郭少雅}

贵州台盘村的乡村篮球赛“出圈”了。一场没有明星、没有赞助的乡村赛事，为什么能吸引全网的目光？因为它让我们看到了乡村文化最生动的模样，也让我们看到了乡村文化自信的回归。

文化自信，源于对自身文化的热爱与认同。台盘村的村民爱篮球，爱了几十年。他们把篮球赛办成自己的节日，把球场当成展示自我的舞台。这份发自内心的热爱，正是文化自信最朴素的表达。当村民骄傲地向外人介绍“我们村的篮球赛”时，我们看到的，是一种由内而外的文化底气。

“村BA”的出圈，也打破了人们对乡村文化的刻板印象。乡村不只有“土”，也有活力；乡村文化不只有“遗产”，也有创造。当互联网把这份乡土热情传递到全国，人们发现：原来乡村文化可以这么鲜活，这么动人。

乡村振兴，既要塑形，也要铸魂。乡村文化自信，就是乡村振兴的“魂”之一。期待更多像“村BA”这样有温度、有活力的乡村文化实践，让广袤的田野上，生长出更多文化自信的花朵。

\clearpage
\section*{技法指南：“塑形 + 铸魂”：精心培育乡村文化生长的土壤}

乡村振兴，既是一场深刻的产业革命，也是一场深刻的文化变革。如果说产业兴旺是乡村振兴的“形”，那么文化振兴就是乡村振兴的“魂”。只有“塑形”与“铸魂”并重，才能让乡村真正“活”起来。

“塑形”，是要筑牢乡村发展的物质基础。完善基础设施、发展特色产业、改善人居环境，让乡村留得住人、引得来客。没有“形”的支撑，“魂”便无处安放。正如费孝通先生在《乡土中国》中所言，乡村社会是一个熟人社会，其秩序建立在人与人之间长期的交往之上。基础设施的改善，让这种交往更加便利，也让乡土文化有了延续的土壤。

“铸魂”，是要激活乡村文化的内生动力。乡村文化不是博物馆里的展品，而是流动在村民日常生活中的活态存在。台盘村的篮球赛、各地的传统节日、手工艺人的匠心，都是乡村文化的鲜活载体。要让村民成为文化的主人，让传统在创新中传承，让文化在交流中生长。

“塑形”与“铸魂”相辅相成。经济强了，文化才有底气；文化兴了，经济才有灵魂。黑格尔曾说：“一个民族有一些关注天空的人，他们才有希望。”乡村振兴同样需要这样“关注天空”的人——他们立足乡土，又放眼未来；他们传承传统，又勇于创新。

从故宫文物修复师的“择一事，终一生”，到《经典咏流传》让古诗词焕发新生，无数实践证明：文化的生命力在于创造性地转化与发展。乡村文化也是如此。当我们用心浇灌，这片古老的土地上，必将生长出更加繁茂的文化之林。

\begin{quoteBox}{技法点评}
本文以“塑形+铸魂”构建全文框架，立意新颖，层次清晰。引用费孝通、黑格尔等名家观点增强说理深度，并结合“故宫修文物”“经典咏流传”等时代案例，做到了古今贯通、虚实结合，是“文化与乡村振兴”类话题的范本写法。
\end{quoteBox}

\clearpage
\section*{课堂演练}
\begin{quote}
\textbf{阅读下面材料，按要求写一篇作文。}

自从贵州台盘村的“村BA”火了以后，其他地方纷纷效仿，也举行“村BA”篮球赛。对此，有人说，走别人走过的路，“人云亦云”，没有意义；也有人说，乡村文化建设不就是应该相互交流和借鉴吗？两者要表达的都是对乡村文化的传承与发扬。

读了以上材料，你有什么感想？请写一篇文章，谈谈自己的看法。

要求：1.角度自选，立意自定，题目自拟，文体不限（除诗歌外）。2.不少于800字。
\end{quote}

\section*{考场范本：借他山之玉，传承民族文化}
\auth{文/广东省佛山市第二中学高二（8）班 林翔}

贵州台盘村的“村BA”火了之后，其他地方纷纷效仿，也举办“村BA”篮球赛。有人担心，这样简单的模仿可能影响本地的文化建设。但事实上，在相互借鉴的过程中融合本地特色，能让乡土文化的形式更多样，底蕴更深厚。村落如此，民族亦然。文化凝聚着集体的智慧，传承和发扬本民族文化是与一个民族生存和发展的重要保障，而学习和借鉴其他民族的文化精华是让本民族文化永葆青春活力的重要途径。

借鉴可以更好地保护自己的文化，找到文化传承的新载体。文化的内涵不会轻易改变，而传播文化的载体却可以灵活选择。中国古籍中本没有“博物馆”这一名词，有的只是“博物”二字，大意就是见多识广，博识多知。为了保存并传播我们的文化，国家积极借鉴国际做法，开设一座座文化博物馆，让一批批文化遗产走进了大众的视野，助力它们的传承与发扬，守住属于本民族的文化精髓。

借鉴可以发现并改良自身的不足，达到“取其精华，去其糟粕”的目的。世界各民族的文化相互借鉴，既能发现对方的优点，又能照出自己的不足。百年前的新文化运动给后人留下深刻的印象，就是因为借鉴了当时西方的先进文化，高举“德先生”与“赛先生”两面大旗，倡导祛除迂腐、落后的旧文化。借鉴他者的优良文化基因，从中发现并改良自身存在的不足，这对于本民族文化的传承与发扬有着举足轻重的作用。

借鉴可以提供新的思路与出路，让文化传承与发扬拥有更广阔的前景。传统文化遇到发展瓶颈，往往不是因为文化本身出了问题，而是缺乏新的发展思路。前不久，贵州省台江县台盘村开展的乡村篮球比赛“村BA”火爆全网，让这个“无名小村”第一次进入全国乃至全球人们的视野，当地村民也利用这个机会宣传了家乡文化。台盘村正是借鉴了“NBA”这一现代篮球文化的表现形式，为传统的乡村文化传承与发展闯出了一条新路。

“独学而无友，则孤陋而寡闻”，《礼记》中这句充满智慧的名言早在两千多年前就阐明了文化交流互鉴的重要性。而在今天，我们更应该意识到，借鉴他民族文化的精髓并为我所用，能够让本民族文化走得更远。
